% Shrinks its content to the text width only when it would be wider than the page.
\newsavebox{\fitbox}
\newcommand{\fitwidth}[1]{\sbox{\fitbox}{#1}%
  \ifdim\wd\fitbox>\linewidth\resizebox{\linewidth}{!}{\usebox{\fitbox}}%
  \else\usebox{\fitbox}\fi}

\chapter{Results}
\label{sec:results}

This chapter presents the empirical evaluation of the proposed context-aware
privacy framework. Section~\ref{sec:dataset} describes the evaluation dataset and ground-truth
annotation protocol. Section~\ref{sec:results_2a} reports classification performance for the
geometry-only (Phase~2a) configuration. Section~\ref{sec:results_2b} extends this to the full
system that additionally incorporates facial landmark cues (Phase~2b), evaluated
first on still frames. Section~\ref{sec:results_clips} reports a second, larger evaluation of Phase~2b
using short video clips with sound. Section~\ref{sec:results_secondpass} reports an offline second pass for
recorded video that preserves confirmed partners and their group for their whole
time on screen. Section~\ref{sec:baseline_comparison} compares all configurations against the
blanket-blurring baseline. Section~\ref{sec:fps} reports computational performance on the
evaluation hardware and provides a projected estimate for the Raspberry~Pi~5
deployment target.

\section{Dataset and Ground-Truth Annotation}
\label{sec:dataset}

The evaluation dataset comprises 21 egocentric video clips that I collected from
publicly available online sources, including YouTube; I did not record any of the
footage myself. All clips are filmed from a first-person viewpoint, with the camera worn
or held at approximately eye level. The clips are divided into two categories:
14 \emph{interaction} clips, in which the camera wearer conducts a face-to-face
conversation with at least one person, and 7 \emph{bystander} clips, in which no
direct interaction takes place and all visible faces belong to people who are
present but not conversing with the wearer.

The two categories differ in setting. The bystander clips show the wearer walking
through busy outdoor public spaces, such as street markets and city streets, by day
and at night, and range from approximately 6 to 78 seconds in length. The
interaction clips are five-minute segments filmed indoors, mostly showing the wearer
playing card or board games with others around a table, with the remainder showing
the wearer shopping in small stores. Source resolutions range from
$1280 \times 720$ to $3840 \times 2160$; all frames were resized to
$1280 \times 720$ for processing. For the Phase~2a evaluation, I manually
labelled 66 face tracks across the 21 clips.

I established the ground-truth labels by watching each clip in full. A track was assigned the label \emph{interaction partner} (1) if the
corresponding person engaged in face-to-face conversation with the camera wearer at
any point during the clip, and \emph{bystander} (0) otherwise. This per-track
binary protocol is consistent with the interaction latching behaviour of the classifier,
which permanently commits a track to the \emph{Interactor} state once the
interaction criteria are met.

Across all 66 tracks, 26 were labelled as interaction partners and 40 as bystanders
(Table~\ref{tab:dataset}).

\begin{table}[htbp]
\begin{singlespace}
\centering
\caption{Evaluation dataset composition and ground-truth track distribution.}
\label{tab:dataset}
\renewcommand{\arraystretch}{1.3}
\begin{tabular}{@{}lc@{}}
\toprule
\textbf{Video clips} & \textbf{Count} \\
\midrule
Interaction clips (5-minute indoor segments) & 14 \\
Bystander clips (6-78~s outdoor walking)    & 7  \\
\textbf{Total clips}                         & \textbf{21} \\
\midrule
\textbf{Labelled face tracks} & \textbf{Count} \\
\midrule
Interaction partner tracks & 26 \\
Bystander tracks           & 40 \\
\textbf{Total tracks}      & \textbf{66} \\
\bottomrule
\end{tabular}
\end{singlespace}
\end{table}

\section{Classification Performance: Phase 2a (Geometry-Only)}
\label{sec:results_2a}

The Phase~2a configuration uses only the two bounding-box geometry cues
(area ratio and center distance) together with temporal hysteresis and latching.
No facial landmark extraction is performed. I evaluated the system against the
66 ground-truth tracks.

Table~\ref{tab:confusion_2a} presents the confusion matrix.

\begin{table}[htbp]
\begin{singlespace}
\centering
\caption{Confusion matrix for the Phase~2a (geometry-only) configuration. Rows give
the actual (human-labelled) class of each face track and columns give the system's
decision.}
\label{tab:confusion_2a}
\renewcommand{\arraystretch}{1.3}
\begin{tabular}{@{}lccc@{}}
\toprule
 & \textbf{Predicted Partner} & \textbf{Predicted Bystander} & \textbf{Total} \\
\midrule
\textbf{Actual partner}   & 22 (TP) & 4 (FN) & 26 \\
\textbf{Actual bystander} & 5 (FP)  & 35 (TN) & 40 \\
\textbf{Total}        & 27      & 39      & 66 \\
\bottomrule
\end{tabular}
\end{singlespace}
\end{table}

The five false positives represent bystanders that the geometry rule incorrectly
classified as partners (privacy leaks), and the four false negatives represent
interaction partners that were incorrectly blurred (utility losses).
Table~\ref{tab:metrics_2a} summarises all derived metrics.

\begin{table}[htbp]
\begin{singlespace}
\centering
\caption{Classification metrics for the Phase~2a (geometry-only) configuration.}
\label{tab:metrics_2a}
\renewcommand{\arraystretch}{1.3}
\begin{tabular}{@{}lcc@{}}
\toprule
\textbf{Metric} & \textbf{Value} & \textbf{Definition} \\
\midrule
Privacy Success Rate (PSR)          & 87.50\% & TN / (TN + FP) \\
Interaction Preservation Rate (IPR) & 84.62\% & TP / (TP + FN) \\
Precision                           & 81.48\% & TP / (TP + FP) \\
Recall                              & 84.62\% & TP / (TP + FN) \\
F$_1$ Score                         & 83.02\% & Harmonic mean of Precision and Recall \\
\bottomrule
\end{tabular}
\end{singlespace}
\end{table}

The Phase~2a system correctly blurred 35 of 40 bystander tracks (PSR = 87.50\%) and
correctly preserved 22 of 26 partner tracks (IPR = 84.62\%). The F$_1$ score of
83.02\% reflects a reasonably balanced trade-off between privacy protection and
interaction preservation.

\section{Classification Performance: Phase 2b (Full System with Landmark Cues)}
\label{sec:results_2b}

The Phase~2b configuration extends Phase~2a by adding the two facial landmark
cues derived from MediaPipe Face Landmarker: head-pose yaw (orientation veto) and
Mouth Aspect Ratio (V-VAD boost). I ran the complete Phase~2b pipeline on all
evaluation clips with the threshold values listed in Table~\ref{tab:params}.

\paragraph{Evaluation protocol.}
The 66-track ground truth of Section~\ref{sec:results_2a} could not be reused directly for Phase~2b.
SORT assigns track identifiers sequentially within each clip, so a single additional
or missing detection early in a clip shifts every subsequent identifier. This
occurred between the Phase~2a and Phase~2b processing runs, and the stored per-track
labels no longer referred to the same individuals. I therefore evaluated Phase~2b
through a blind, multi-annotator manual annotation study.

Ten evaluation frames were drawn automatically, one from each of the first five
bystander clips and the first five interaction clips. Within each clip, the frame
with the most clearly visible faces was selected (bounding box of at least
50~pixels per side and detection confidence of at least 0.50, excluding the first
30 frames while the tracker initialises). Each face in the selected frames was
marked with a plain numbered badge, with no colour or label indicating the system's
decision. Four annotators (myself and three independent annotators)
labelled every face as \emph{interaction partner} or \emph{bystander}, or left it
blank if undecidable; 114 of the 116 possible labels were provided. The final label
of each face was obtained by majority vote, with ties resolved as \emph{bystander}
(the privacy-conservative choice); one tie occurred. Each annotated face corresponds
to a SORT track, and, consistent with the latching protocol of Section~\ref{sec:dataset}, the
system's prediction for a face is \emph{partner} if its track was classified as
\textsc{Interactor} at any point in the clip.

I measured inter-annotator agreement with Cohen's~$\kappa$ for each of the six
annotator pairs. Pairwise $\kappa$ ranged from 0.307 to 0.724 (mean 0.587), with
five of the six pairs showing moderate or substantial agreement. The majority vote
yielded 29 annotated faces: 7 interaction partners and 22 bystanders
(Table~\ref{tab:sample_2b}).

\begin{table}[htbp]
\begin{singlespace}
\centering
\caption{Composition of the human-annotated Phase~2b evaluation sample.}
\label{tab:sample_2b}
\renewcommand{\arraystretch}{1.3}
\begin{tabular}{@{}lcccc@{}}
\toprule
\textbf{Clip category} & \textbf{Frames} & \textbf{Faces} & \textbf{Partners} & \textbf{Bystanders} \\
\midrule
Interaction & 5  & 15 & 7 & 8  \\
Bystander   & 5  & 14 & 0 & 14 \\
\midrule
\textbf{Total} & \textbf{10} & \textbf{29} & \textbf{7} & \textbf{22} \\
\bottomrule
\end{tabular}
\end{singlespace}
\end{table}

\paragraph{Results.}
Table~\ref{tab:confusion_2b} presents the confusion matrix for Phase~2b against the
majority-vote labels.

\begin{table}[htbp]
\begin{singlespace}
\centering
\caption{Confusion matrix for the Phase~2b (full system) configuration on the
human-annotated sample. Rows give the actual (human-labelled) class of each face and
columns give the system's decision.}
\label{tab:confusion_2b}
\renewcommand{\arraystretch}{1.3}
\begin{tabular}{@{}lccc@{}}
\toprule
 & \textbf{Predicted Partner} & \textbf{Predicted Bystander} & \textbf{Total} \\
\midrule
\textbf{Actual partner}   & 4 (TP) & 3 (FN)  & 7 \\
\textbf{Actual bystander} & 0 (FP) & 22 (TN) & 22 \\
\textbf{Total}        & 4      & 25      & 29 \\
\bottomrule
\end{tabular}
\end{singlespace}
\end{table}

\begin{table}[htbp]
\begin{singlespace}
\centering
\caption{Classification metrics for the Phase~2a and Phase~2b configurations.}
\label{tab:metrics_2b}
\renewcommand{\arraystretch}{1.3}
\fitwidth{\begin{tabular}{@{}lcc@{}}
\toprule
\textbf{Metric} & \textbf{Phase 2a} & \textbf{Phase 2b (Full System)} \\
\midrule
Evaluation set                      & 66 tracks, 1 annotator & 29 faces, 4 annotators \\
\midrule
Privacy Success Rate (PSR)          & 87.50\% & 100.00\% \\
Interaction Preservation Rate (IPR) & 84.62\% & 57.14\% \\
Precision                           & 81.48\% & 100.00\% \\
Recall                              & 84.62\% & 57.14\% \\
F$_1$ Score                         & 83.02\% & 72.73\% \\
\bottomrule
\end{tabular}}
\end{singlespace}
\end{table}

Phase~2b produced no false positives: all 22 bystanders in the sample were blurred
(PSR = 100\%), including the 8 bystanders who appeared in interaction clips
alongside an active partner. Of the 7 interaction partners, 4 were correctly
preserved and 3 were blurred (IPR = 57.14\%). Broken down by clip category, all 14
faces in bystander clips were correctly blurred, and in interaction clips the
system preserved 4 of 7 partners while blurring all 8 bystanders.

Because the two configurations were evaluated on different samples, the
comparison in Table~\ref{tab:metrics_2b} is indicative rather than conclusive. The
Phase~2b sample is small: with only seven partners, each misclassified partner
changes IPR by 14.3 percentage points. The 95\% Wilson confidence intervals are
85.1-100\% for Phase~2b PSR and 25.0-84.2\% for Phase~2b IPR, compared with
73.9-94.5\% (PSR) and 66.5-93.9\% (IPR) for Phase~2a; the intervals overlap for
both metrics. The causes of the three Phase~2b false negatives are analysed in
Section~\ref{sec:disc_failures}.

\section{Phase 2b: Video-Clip Evaluation}
\label{sec:results_clips}

The still-frame study had two weaknesses. Its sample was small (seven partners),
and annotators had to judge interaction from a single image, which explains its
moderate agreement. I therefore ran a second, larger blind evaluation of the same
Phase~2b system using short video clips with sound.

\paragraph{Evaluation protocol.}
I sampled face tracks from all 21 evaluation clips without reference to the
system's decisions: in each interaction clip, the four longest tracks plus two
randomly chosen tracks, and in each bystander clip, the four longest tracks. A track
was eligible if it was visible for at least 60 frames (about two seconds) and its
median face size was at least 32~pixels. This gave 104 clips (83 from interaction
videos and 21 from bystander videos). Each clip lasts seven seconds: a one-second
still frame showing which face to watch, followed by six seconds of the original,
unblurred video with its original sound. A single yellow box marked the target face,
and the clips were shown in a shuffled order with no indication of their source
category or of the system's decision.

For this study I defined an \emph{interaction partner} as any member of the camera
wearer's group or conversation. This includes a person talking to the wearer and
also a person who is with the wearer but not speaking at that moment, such as a
friend at the same table or playing the same game. A \emph{bystander} is anyone not
in the wearer's group, such as a stranger walking past or a person in a different
group. This is broader than the still-frame instruction (actively talking to the
wearer), and reflects the social reality that the members of the wearer's own group
are not the people whose privacy is at risk. Five annotators (myself and four
independent annotators) labelled each clip through a simple web page. The final
label was obtained by majority vote; with five annotators, no ties occurred. In two clips the boxed region was
not a person at all but a false detection by the face detector: a face printed on a
magazine cover and a product package on a shop shelf. Most annotators left these
clips unlabelled and I excluded them, leaving 102 labelled clips: 69 partners and 33
bystanders.

The system's decision for each clip is the decision visible in its output video: a
face counts as preserved if it was left unblurred in at least half of the frames in
which it was detected during the clip. I verified these decisions directly against
the rendered output videos, and all 104 agreed with the decisions recorded by the
pipeline.

\paragraph{Agreement.}
Agreement between annotators was almost perfect (Table~\ref{tab:clip_kappa}). The
pairwise Cohen's~$\kappa$ ranged from 0.818 to 0.957 and Fleiss'~$\kappa$ across all
five annotators was 0.887, compared with a mean pairwise $\kappa$ of 0.587 in the
still-frame study. Seeing movement and hearing speech therefore made the labelling
task far more reliable.

\begin{table}[htbp]
\begin{singlespace}
\centering
\caption{Inter-annotator agreement in the video-clip evaluation. Annotator A is
myself.}
\label{tab:clip_kappa}
\renewcommand{\arraystretch}{1.3}
\begin{tabular}{@{}lcc@{}}
\toprule
\textbf{Annotator pair} & \textbf{Cohen's $\kappa$} & \textbf{Raw agreement} \\
\midrule
A and B & 0.957 & 98.0\% \\
A and C & 0.913 & 96.1\% \\
A and D & 0.818 & 92.1\% \\
A and E & 0.822 & 92.2\% \\
B and C & 0.956 & 98.0\% \\
B and D & 0.861 & 94.1\% \\
B and E & 0.864 & 94.1\% \\
C and D & 0.905 & 96.0\% \\
C and E & 0.907 & 96.1\% \\
D and E & 0.852 & 94.1\% \\
\midrule
\textbf{All five (Fleiss' $\kappa$)} & \textbf{0.887} & \\
\bottomrule
\end{tabular}
\end{singlespace}
\end{table}

\paragraph{Results.}
Table~\ref{tab:confusion_clips} gives the confusion matrix and
Table~\ref{tab:metrics_clips} the metrics, overall and by clip category.

\begin{table}[htbp]
\begin{singlespace}
\centering
\caption{Confusion matrix for Phase~2b in the video-clip evaluation. Rows give the
actual (human-labelled) class of each face and columns give the system's decision.}
\label{tab:confusion_clips}
\renewcommand{\arraystretch}{1.3}
\begin{tabular}{@{}lccc@{}}
\toprule
 & \textbf{Predicted Partner} & \textbf{Predicted Bystander} & \textbf{Total} \\
\midrule
\textbf{Actual partner}   & 39 (TP) & 30 (FN) & 69 \\
\textbf{Actual bystander} & 1 (FP)  & 32 (TN) & 33 \\
\textbf{Total}            & 40      & 62      & 102 \\
\bottomrule
\end{tabular}
\end{singlespace}
\end{table}

\begin{table}[htbp]
\begin{singlespace}
\centering
\caption{Phase~2b metrics in the video-clip evaluation, with 95\% Wilson confidence
intervals. IPR is not defined for bystander clips, which contain no partners.}
\label{tab:metrics_clips}
\renewcommand{\arraystretch}{1.3}
\fitwidth{\begin{tabular}{@{}lccc@{}}
\toprule
\textbf{Metric} & \textbf{All clips} & \textbf{Interaction clips} & \textbf{Bystander clips} \\
\midrule
Faces evaluated & 102 & 81 & 21 \\
PSR       & 96.97\% (84.7-99.5) & 100.00\% (75.7-100.0) & 95.24\% (77.3-99.2) \\
IPR       & 56.52\% (44.8-67.6) & 56.52\% (44.8-67.6) & not defined \\
Precision & 97.50\% & 100.00\% & not defined \\
F$_1$     & 71.56\% & 72.22\% & not defined \\
\bottomrule
\end{tabular}}
\end{singlespace}
\end{table}

The system blurred 32 of 33 bystanders (PSR = 96.97\%), including 20 of the 21 faces
in bystander clips and all 12 bystanders who appeared in interaction clips alongside
the wearer's group. When the system left a face unblurred, it was correct 97.5\% of
the time. It preserved 39 of 69 partners (IPR = 56.52\%). These figures closely match
the still-frame study (PSR 100\%, IPR 57.14\%), but the IPR confidence interval is now
about half as wide (44.8-67.6\% against 25.0-84.2\%), so the video-clip study confirms
the earlier finding with much stronger evidence. Of the 30 partners that were
blurred, 25 had a median face size below the 1.2\% area threshold, and 10 were
confirmed as partners by the system later in the same video, after the six-second
clip had ended. These
failures are analysed in Section~\ref{sec:disc_failures}.

\section{Offline Second Pass for Recorded Video}
\label{sec:results_secondpass}

The interaction latch already ensures that a confirmed partner is never blurred
again. However, it works only forward in time: a partner is blurred until the
system has gathered enough evidence to confirm them. When the video is recorded and
processed after capture, the whole of each track is known in advance, so this
limitation can be removed. I implemented a second processing pass over the recorded
decisions, which applies two rules.

\begin{enumerate}
  \item \textbf{Retroactive latch.} A track that is confirmed as a partner at any
        point is preserved for its whole time on screen, including the frames before
        confirmation. This applies the principle that once a person is an interaction
        partner, they remain one.
  \item \textbf{Friend group.} A track that stays in view for at least ten seconds,
        and is on screen together with a confirmed partner for at least half of that
        time, is treated as a member of the wearer's group and preserved as well.
        This covers friends who sit with the wearer but do not speak, while strangers
        who pass by are rarely in view for that long. Requiring the presence of a
        confirmed partner prevents the rule from preserving strangers who simply stay
        in one place, such as a shop assistant.
\end{enumerate}

Neither rule requires additional model inference, since both operate on the
decisions the pipeline has already made. Table~\ref{tab:secondpass} reports their
effect on the 102 clips of the video-clip evaluation.

\begin{table}[htbp]
\begin{singlespace}
\centering
\caption{Effect of the offline second pass on the video-clip evaluation (102 clips,
95\% Wilson confidence intervals).}
\label{tab:secondpass}
\renewcommand{\arraystretch}{1.3}
\fitwidth{\begin{tabular}{@{}lccc@{}}
\toprule
\textbf{Metric} & \textbf{Live system} & \textbf{+ Retroactive latch} & \textbf{+ Friend group} \\
\midrule
TP / FN / FP / TN & 39 / 30 / 1 / 32 & 49 / 20 / 2 / 31 & 54 / 15 / 2 / 31 \\
PSR       & 96.97\% (84.7-99.5) & 93.94\% (80.4-98.3) & 93.94\% (80.4-98.3) \\
IPR       & 56.52\% (44.8-67.6) & 71.01\% (59.4-80.4) & 78.26\% (67.2-86.4) \\
Precision & 97.50\% & 96.08\% & 96.43\% \\
F$_1$     & 71.56\% & 81.67\% & 86.40\% \\
\bottomrule
\end{tabular}}
\end{singlespace}
\end{table}

The retroactive latch recovered 10 partners at the cost of one additional
bystander being preserved. The friend-group rule recovered a further 5 partners
without any additional privacy leak. Together, the two rules raised IPR from 56.52\%
to 78.26\% and F$_1$ from 71.56\% to 86.40\%, while PSR fell from 96.97\% to 93.94\%.
To check that the ten-second limit was not simply fitted to these clips, I repeated
the evaluation with limits from 10 to 60~seconds. Every limit in this range gave the
same PSR (93.94\%) and an IPR between 72.46\% and 78.26\%, so the result is not
sensitive to the exact value. Only a much shorter limit of five seconds admitted an
additional bystander.

Both rules were designed after examining the errors of the video-clip evaluation and
are evaluated on the same labels, so these figures should be read as an indication
of the benefit of the second pass rather than as an independent test. Both rules
also require the whole recording and therefore apply to recorded video, not to
live streaming.

\section{Comparison with Blanket-Blurring Baseline}
\label{sec:baseline_comparison}

The blanket-blurring baseline applies Gaussian blur to every detected face in every
frame, making no use of interaction cues. By construction, this baseline achieves a
PSR of 100\% (every bystander is blurred) and an IPR of 0\% (every interaction
partner is also blurred). The proposed system represents a qualitative departure
from the baseline: Phase~2a sacrifices a small amount of PSR to recover interaction
utility, while Phase~2b stays close to the baseline's PSR and still recovers a
large part of the interaction utility (Table~\ref{tab:vs_baseline}).

\begin{table}[htbp]
\begin{singlespace}
\centering
\caption{Comparison of the proposed system against the blanket-blurring baseline.
Phase~2a was evaluated on the 66-track ground truth, Phase~2b on the 29-face
still-frame sample (Section~\ref{sec:results_2b}) and on the 102-clip video sample
(Sections~\ref{sec:results_clips} and~\ref{sec:results_secondpass}).}
\label{tab:vs_baseline}
\renewcommand{\arraystretch}{1.3}
\begin{tabular}{@{}lccc@{}}
\toprule
\textbf{Configuration} & \textbf{PSR} & \textbf{IPR} & \textbf{F$_1$} \\
\midrule
Blanket-blur baseline      & 100\%   & 0\%     & N/A \\
Phase~2a (geometry)        & 87.50\% & 84.62\% & 83.02\% \\
Phase~2b, still frames     & 100.00\% & 57.14\% & 72.73\% \\
Phase~2b, video clips      & 96.97\% & 56.52\% & 71.56\% \\
Phase~2b + offline second pass & 93.94\% & 78.26\% & 86.40\% \\
\bottomrule
\end{tabular}
\end{singlespace}
\end{table}

The baseline's perfect PSR is achieved at the cost of complete utility destruction:
the camera wearer's interaction partners are indistinguishable from strangers in the
output stream. The proposed Phase~2a system recovers 84.62\% of interaction partners
(IPR) while correctly anonymizing 87.50\% of bystanders (PSR), at no additional
hardware cost relative to the baseline, since both configurations share the same
YOLOv8n-face detector and SORT tracker. In the video-clip evaluation, the Phase~2b
system blurred 96.97\% of bystanders while preserving 39 of 69 partners that the
baseline would have blurred, and with the offline second pass it preserved 54 of
69 partners (78.26\%) while still blurring 93.94\% of bystanders.

\section{Computational Performance}
\label{sec:fps}

I measured pipeline latency by re-running the full processing loop on a laptop
CPU (Intel Core i-series, no GPU or NPU acceleration enabled) and timing each stage
independently with a high-resolution wall-clock timer over a 300-frame clip. Table~\ref{tab:fps}
reports the per-stage mean latency for the Phase~2a configuration.

\begin{table}[htbp]
\begin{singlespace}
\centering
\caption{Per-stage mean latency (Phase~2a, 300-frame clip, laptop CPU).}
\label{tab:fps}
\renewcommand{\arraystretch}{1.3}
\begin{tabular}{@{}lcc@{}}
\toprule
\textbf{Stage} & \textbf{Mean latency (ms)} & \textbf{Fraction of total} \\
\midrule
Frame read          & 10.0  & 11.6\% \\
Face detection (YOLOv8n-face) & 55.1 & 64.1\% \\
Tracking (SORT)     & 20.1  & 23.4\% \\
Classification      & 0.02  & $<$0.1\% \\
Selective anonymization & 0.87  & 1.0\% \\
\midrule
\textbf{Total per frame} & \textbf{86.0} & \textbf{100\%} \\
\textbf{Throughput}      & \multicolumn{2}{c}{\textbf{11.62 FPS}} \\
\bottomrule
\end{tabular}
\end{singlespace}
\end{table}

Face detection with YOLOv8n-face dominates the pipeline at 55.1~ms per frame
(64.1\% of total latency). The classification stage is negligible at 0.02~ms,
confirming that the heuristic decision layer imposes no measurable overhead relative
to the underlying detection and tracking stages.

\paragraph{Raspberry Pi 5 projection.}
I did not benchmark directly on a Raspberry Pi 5, as the target hardware was not
available to me at the time of evaluation. Following the methodology described in
Section~3.8, a conservative hardware scaling factor of 3.5$\times$ is applied to
project performance on the ARM Cortex-A76 processor of the Raspberry Pi 5. This
factor is derived from published YOLO inference benchmarks comparing Intel Core-class
and ARM Cortex-A76 class processors on equivalent workloads. Applying this factor to
the measured 86.0~ms per-frame latency yields an estimated throughput of
approximately \textbf{3.3 FPS} on the Raspberry Pi 5 for the Phase~2a configuration.

\paragraph{Phase 2b landmark overhead.}
I benchmarked the landmark stage separately by timing every MediaPipe Face
Landmarker call within the same 300-frame loop, on the same laptop CPU under mains
power. Across two runs, landmark extraction measured \textbf{20.1~ms per face crop}
on average (21.1 and 19.1~ms); the benchmark clip averaged 0.99 tracked faces per
frame, adding 19.9~ms per frame. This is consistent with the 15-25~ms per-face
estimate anticipated in the methodology. Because detection latency varied between
runs owing to background CPU load, the Phase~2b figure is obtained by adding the
measured landmark cost to the Phase~2a baseline: $86.0 + 19.9 \approx 106$~ms per
frame, or approximately \textbf{9.4 FPS} on the laptop CPU (the fastest directly
measured Phase~2b run reached 9.9~FPS, in agreement). Applying the same
3.5$\times$ scaling factor gives approximately \textbf{2.7 FPS} on the Raspberry
Pi 5. Since landmark extraction runs once per tracked face, this overhead grows
linearly with the number of faces in view.

\paragraph{Battery power.}
Wearable devices always run on battery, and a laptop on battery limits its processor
speed to save energy. I therefore also benchmarked the pipeline with the laptop
unplugged. On battery, landmark extraction took
33.5~ms per face, and the complete Phase~2b pipeline measured 163.8~ms per frame
(6.1~FPS); two Phase~2a runs measured 9.1 and 7.5~FPS. These runs were also
affected by rising processor temperature over consecutive runs, so they represent a
conservative, worst-case condition rather than a precise figure.
Table~\ref{tab:power} summarises both conditions.

\begin{table}[htbp]
\begin{singlespace}
\centering
\caption{Measured laptop CPU throughput under mains and battery power, with the
projected Raspberry Pi 5 throughput (3.5$\times$ scaling factor). Battery figures are
single runs affected by rising processor temperature.}
\label{tab:power}
\renewcommand{\arraystretch}{1.3}
\fitwidth{\begin{tabular}{@{}lcccc@{}}
\toprule
 & \multicolumn{2}{c}{\textbf{Mains power}} & \multicolumn{2}{c}{\textbf{Battery power}} \\
\textbf{Configuration} & \textbf{Laptop} & \textbf{Pi 5 (projected)} & \textbf{Laptop} & \textbf{Pi 5 (projected)} \\
\midrule
Phase~2a & 11.6~FPS & 3.3~FPS & 7.5-9.1~FPS & 2.1-2.6~FPS \\
Phase~2b & 9.4~FPS  & 2.7~FPS & 6.1~FPS     & 1.7~FPS \\
Landmarks per face & 20.1~ms & & 33.5~ms & \\
\bottomrule
\end{tabular}}
\end{singlespace}
\end{table}

In summary, the complete pipeline runs entirely on a CPU, without a GPU, at
11.62~FPS (Phase~2a) and approximately 9.4~FPS (Phase~2b) on a laptop under mains
power, and at 6.1-9.1~FPS on battery, with projected Raspberry Pi 5 throughput of
1.7-3.3~FPS across both conditions. The interaction classifier
itself adds less than 0.1\% to the processing time, so the selective decision comes
at almost no extra cost over blanket blurring. Section~\ref{sec:disc_fps} discusses optimisations
that can raise throughput further on edge hardware.
